\section{What can neural measures reveal about how receivers assess a sound?}
\label{sec:q3}

% Part II — Assessment and behavioral response (Q3). Keep deliberately short.
% See docs/OUTLINE.md — Q3.
% Draft 1. Style model: Watabe-Uchida (agents/STYLE.md).

Neural measures are useful to this review to the extent that they reveal which computation a receiver performs on a sound: detecting its acoustic features, assigning it a biological interpretation, selecting an action, or updating expectations. A neural response is not a direct readout of perceived threat or of the behavioral decision that follows, and the methods available differ in what they permit one to infer. We therefore restrict this section to contrasts in which a neural measure dissociates from behavior or from a simple account of stimulus value.

The first such contrast shows that discrimination can occur without overt behavior and need not be specific to natural calls. In naïve rats, 22-kHz calls did not elicit freezing, yet amygdala neurons responded differently to 22-kHz and 50-kHz calls \citep{ParsanaEtAl2012a}. Tones matched in frequency reproduced most of this difference, so the difference largely reflects a code at the level of acoustic detection rather than a representation of call meaning.
% VERIFY: EilamEtAl1999 and Mateo1996 (not full-text checked): endocrine response without behavior; heart-rate discrimination preceding behavior.
Physiological measures point the same way: owl calls can raise glucocorticoids without behavioral change \citep{EilamEtAl1999}, and juvenile ground squirrels discriminate alarm calls by heart rate before they do so behaviorally \citep{Mateo1996} (\secref{sec:q2}, \secref{sec:q4}). Absence of a behavioral response is not absence of assessment, but presence of a neural difference is not evidence that the receiver assigned the sounds different meanings.

Dopamine recordings test whether a single signal tracks the value of a sound, and the answer depends on region and preparation. Appetitive 50-kHz calls transiently increased dopamine in the rat nucleus accumbens, whereas 22-kHz calls did not \citep{WilluhnEtAl2014}. In head-fixed mice exposed to noise bursts between 45 and 105~dB, accumbens dopamine fell in all animals and the decrease grew with intensity; six of fifteen mice also showed an increase at the onset of the first burst, consistent with a novelty response, while dopamine in the tail of the striatum rose \citep{FuruyamaEtAl2026}. In freely moving rats, accumbens core dopamine also fell during aversive noise and during a cue that predicted it, tracking the occurrence, prediction, and duration of the noise, but it did not scale with intensity between 70 and 96~dB although avoidance behavior did \citep{GoedhoopEtAl2022}. The two results agree on the direction of the accumbens response and disagree on intensity scaling. Furuyama and colleagues note that head fixation removes the option to escape, which could change how the same stimulus is valued; species and recording method also differ. Neither study supports treating accumbens dopamine as a scalar measure of threat, and the opposite responses of the accumbens and tail of the striatum to the same noise argue against treating striatal dopamine as a single signal.

Activity-dependent markers in a field species caution against equating threat with greater activation. In Richardson's ground squirrels tested in a laboratory chamber, playback of alarm chirps reduced c-Fos expression in the medial amygdala and superior colliculus relative to controls that heard no call, and most other regions did not differ \citep{FreemanEtAl2020}. Because behavior was constrained in the chamber, the result cannot be related to the vigilance and escape responses these calls evoke in the field.
% VERIFY: SadanandaEtAl2008 (not full-text checked): region-specific activation by 50-kHz versus 22-kHz playback.
In laboratory rats, 50-kHz and 22-kHz playback activated partly distinct sets of regions \citep{SadanandaEtAl2008}, but immediate-early gene expression integrates activity over many minutes and cannot distinguish detection, interpretation, and action selection.

These results define what neural measures can and cannot contribute to the questions that follow. If receivers track the reliability of a caller (\secref{sec:q5}), a neural correlate of reliability should change with the history of a specific caller and recover when the caller changes, rather than declining with repetition alone. If receivers form expectations about timing (\secref{sec:q7}), responses should depend on whether a sound was predicted, as accumbens dopamine recovery after noise does \citep{GoedhoopEtAl2022}. These are hypotheses about expected signatures, not findings. Testing them is limited by a gap in the current literature: dopamine responses have been measured mainly to noise, to appetitive 50-kHz calls, or to non-acoustic aversive stimuli, and apart from the 22-kHz calls of \citet{WilluhnEtAl2014}, we found no recordings during natural alarm or predator calls.

\tabref{tab:summary-q3} consolidates these claims for later condensation or figure drafting.

% Intermediate summary matrix for Q3. May later be condensed, moved to an appendix, or redrawn as a figure.
\begin{table}[tbp]
\centering
\footnotesize
\singlespacing
\caption{Working summary for Q3: what neural measures can and cannot show. Each row separates a warranted inference from a common over-reading.}
\label{tab:summary-q3}
\begin{tabular}{@{}p{1.45in}p{2.4in}p{2.4in}@{}}
\toprule
\textbf{Observation} & \textbf{Warranted inference} & \textbf{Not warranted} \\
\midrule
Amygdala differs for 22- vs 50-kHz calls; matched tones reproduce most of it & Acoustic detection / frequency code & Call-specific meaning; threat encoding \\
\addlinespace[2pt]
50-kHz raises, 22-kHz does not raise, NAc dopamine & Valence-sensitive transient & Unitary reward or fear signal \\
\addlinespace[2pt]
Noise lowers NAc DA, raises tail-striatum DA; intensity scaling depends on preparation & Region- and context-specific responses & Scalar threat code in NAc \\
\addlinespace[2pt]
Chirps reduce c-Fos in medial amygdala and superior colliculus & Threat $\neq$ more activation & Mapping to field escape/vigilance \\
\addlinespace[2pt]
IEG patterns differ by call type & Regional engagement over minutes & Detection vs interpretation vs action \\
\addlinespace[2pt]
DA to ecological threat calls & Sparse / largely untested & --- \\
\bottomrule
\end{tabular}
\end{table}

